% Copyright (c) 2026 Geovana Neves. Licensed under CC BY 4.0 (https://creativecommons.org/licenses/by/4.0/). See LICENSE-AND-AUTHORSHIP.md.
%
\section{\code{phase\_locked}}

\begin{frame}{A mean at a fixed azimuth across revolutions}
\begin{enumerate}\small
  \item Take the last \code{last\_revolutions\_avg} revolutions.
  \item At each azimuth, average the samples at that same position across
        those revolutions.
  \item Tabulate the result by azimuthal position, 0 to 360 degrees.
\end{enumerate}
\begin{takeaway}
With three revolutions, three datapoints enter each azimuth's mean.
The row is an azimuthal position, not a passage, revolution or blade.
\end{takeaway}
\end{frame}

\begin{frame}{Blade and alias means in one file}
\begin{itemize}
  \item Do the fixed-azimuth average for each blade on its own, about that
        rotor's \code{SMRP}.
  \item Then do it for each alias, about the global \code{MRP}.
  \item Put all of it in one file with columns suffixed
        \code{\_\{alias or blade\_name\}\_\{SMRP or MRP\}}.
\end{itemize}
\end{frame}

\begin{frame}{The azimuthal table's leading columns}
\small \code{probes/<point>\_phase\_locked[\_<ALIAS>].csv} leads with:
\begin{enumerate}
  \item \code{REDUCTION}, \code{ROTOR}, \code{AZIMUTH}, \code{STEP},
        \code{REVOLUTIONS}.
  \item The steps the revolutions span:
        \code{FIRST\_STEP}, \code{LAST\_STEP}, \code{STEPS}.
  \item The condition block and the moment point.
  \item The plotted columns, under the names the export prints.
\end{enumerate}
\end{frame}

\begin{frame}{Rows use the rotor's last revolution}
\begin{itemize}
  \item There is one row per solver step of the rotor's last revolution,
        sorted by azimuth from 0 towards 360.
  \item \code{AZIMUTH} is where blade one is, using the sections-table formula.
  \item \code{STEP} is the step of the last revolution that azimuth falls on.
  \item \code{REVOLUTIONS} is the number of samples entering the mean.
        It equals \code{last\_revolutions\_avg} when that is a whole number.
\end{itemize}
\end{frame}

\begin{frame}{Blade columns align by each blade's own azimuth}
\begin{itemize}\small
  \item A column ending in a blade family of the rotor is sampled where
        \textbf{that blade} is at the row's azimuth.
  \item The blade sits its position times $360/\mathrm{blades}$ after blade one.
        This offset makes two blades line up azimuth for azimuth.
  \item Every other column, including a rotor total or the aircraft, is
        tabulated by blade one's azimuth.
\end{itemize}
\end{frame}

\begin{frame}{Azimuthal suffixes come from the plots}
\begin{itemize}\small
  \item A plot is \code{<parameter>\_<group>}. The pproc names the group.
  \item The suffix \code{\_\{alias or blade\_name\}\_\{SMRP or MRP\}}
        is what a group named \code{PUSHER\_SMRP} or cut per blade prints.
  \item The package does not rename these columns.
\end{itemize}
\end{frame}

\begin{frame}{Samples between history steps are read linearly}
\begin{itemize}
  \item One revolution earlier falls a whole number of steps earlier only
        when \code{steps\_per\_revolution} is whole.
  \item A blade offset falls on a whole step only when the steps per
        revolution divide by the blade count.
  \item Where both hold, every sample is a history row and nothing is
        interpolated. Between steps, the history is read linearly.
\end{itemize}
\end{frame}

\begin{frame}{The azimuthal mean needs a rotor clock and depth}
\begin{itemize}
  \item Without blade one's datum or the rotor's signed speed, an azimuth
        cannot be stated. The table is a named skip.
  \item Declare the rotor in the reference and have the matrix row cite it.
  \item A depth below one revolution is a named skip too.
\end{itemize}
\end{frame}

\begin{frame}[fragile]{The pproc declares the gate and averaging depth}
\begin{lstlisting}[style=ftsbase]
[phase_locked]
min_revolutions      = 4.0
last_revolutions_avg = 2.0
\end{lstlisting}
\begin{itemize}\small
  \item \code{min\_revolutions} is the minimum \textbf{total} revolutions
        required for the table to exist. Equal or greater generates it.
  \item Compare the revolutions the matrix row states it turns,
        not the length of the exported window.
  \item \code{last\_revolutions\_avg} is the depth of each azimuthal mean
        and may not exceed \code{min\_revolutions}.
\end{itemize}
\end{frame}

\begin{frame}{The gate and depth can change at post}
\begin{itemize}
  \item The declared \code{[phase\_locked]} table binds.
  \item \code{pyfs-matrix post} reads the current pproc again.
  \item Editing the gate or depth needs no solver re-run.
\end{itemize}
\end{frame}

\begin{frame}{A short run loses only the azimuthal reduction}
\begin{itemize}
  \item A row turning fewer total revolutions than the minimum skips this
        reduction.
  \item Its POLAR, per-blade table and every other product are written as usual.
\end{itemize}
\begin{takeaway}
Missing the threshold skips the reduction; it does not refuse the run's products.
\end{takeaway}
\end{frame}

\begin{frame}{An undeclared table keeps the passage series}
\begin{itemize}
  \item A pproc without \code{[phase\_locked]} still gets the averaging
        window cut into consecutive blade passages, one row per passage.
  \item This is the passage series written under this name.
  \item Declaring the table selects the fixed-azimuth mean instead.
        Leaving it absent preserves the passage series.
\end{itemize}
\begin{takeaway}
Absent is not zero: there is no minimum gate on the undeclared passage series.
\end{takeaway}
\end{frame}
